\documentclass[12pt]{article}
\usepackage{amsmath, amssymb, amsthm}
\usepackage{enumitem}
\usepackage[margin=1in]{geometry}

\begin{document}

\section*{Selected IMO Problem}

\begin{enumerate}
    \setcounter{enumi}{3}
    \item For positive real numbers $a, b, c$ with $a + b + c = abc$, show that
    \[
        \frac{1}{\sqrt{1+a^2}} + \frac{1}{\sqrt{1+b^2}} + \frac{1}{\sqrt{1+c^2}} \le \frac{3}{2},
    \]
    and determine when equality occurs.
\end{enumerate}

\section*{English-Language Solution}
\textbf{Proof by:} Trigonometric Substitution \& Jensen's Inequality

\begin{enumerate}
    \item Test $a, b, c$ values when they all equal one another: $(a = b = c \implies 3a = a^3) \implies a = \sqrt{3}$.
    \begin{enumerate}
        \item Plug in $a = b = c = \sqrt{3}$ to get $3\left(\frac{1}{\sqrt{4}}\right) = \frac{3}{2}$.
        \begin{enumerate}
            \item This shows that $\frac{3}{2}$ is possible when all the variables equal $\sqrt{3}$.
        \end{enumerate}
    \end{enumerate}

    \item Using the tangent identity for a triangle (conditional trigonometric identity):
    \begin{enumerate}
        \item $\tan(A) + \tan(B) + \tan(C) = \tan(A)\tan(B)\tan(C)$
        \begin{enumerate}
            \item Here, the angles must add up to $180^\circ$ ($\pi$ radians).
            \item Which can be proven by the Tangent Sum Identity: $\tan(A+B) = \frac{\tan(A) + \tan(B)}{1 - \tan(A)\tan(B)} = -\tan(C)$
            \item Using this identity and the nature of the tan function's period, and odd symmetry, we proved that the 3 angles must be equal to 180 to be true for the Tangent Sum Identity. 
            \item \textbf{Formal Proof Segment:}
            \begin{itemize}
                \item \textbf{Apply Odd Symmetry:} Since the tangent function is an odd function (meaning $-\tan(x) = \tan(-x)$ for all $x$ in its domain), we can rewrite the right-hand side of $\tan(A + B) = -\tan(C)$ as:
                \[
                    \tan(A + B) = \tan(-C)
                \]
                \item \textbf{Apply Periodicity:} The tangent function is strictly periodic with fundamental period $\pi$ (or $180^\circ$). Therefore, $\tan(\alpha) = \tan(\beta)$ if and only if $\alpha - \beta = k\pi$ for some integer $k \in \mathbb{Z}$. Applying this to $\alpha = A + B$ and $\beta = -C$ yields:
                \[
                    (A + B) - (-C) = k\pi \implies A + B + C = k\pi, \quad \text{for some } k \in \mathbb{Z}
                \]
                \item \textbf{Bound the Sum to Isolate $k$:} Since $A, B, C \in \left(0, \frac{\pi}{2}\right)$, summing these strict inequalities gives:
                \[
                    0 < A + B + C < \frac{3\pi}{2}
                \]
                Substituting $A + B + C = k\pi$ into the inequality:
                \[
                    0 < k\pi < \frac{3\pi}{2}
                \]
                Dividing through by $\pi$ gives:
                \[
                    0 < k < 1.5
                \]
                \item \textbf{Conclude:} Because $k$ must be an integer ($k \in \mathbb{Z}$), the unique solution is $k = 1$. Thus:
                \[
                    A + B + C = \pi \quad (\text{or } 180^\circ)
                \]
            \end{itemize}
        \end{enumerate}
    \end{enumerate}

    \item Set $a = \tan(A)$, $b = \tan(B)$, $c = \tan(C)$.
    \begin{enumerate}
        \item Simplify using a right triangle to get $a = \text{opp.}$ and $1 = \text{adj.}$
        \begin{enumerate}
            \item You can rewrite this to show $\cos(A) = \frac{\text{adj.}}{\text{hyp.}} = \frac{1}{\sqrt{1 + a^2}}$.
        \end{enumerate}
        \item Apply the same substitution to the rest of the expression, rewriting the goal as $\cos(A) + \cos(B) + \cos(C) \le \frac{3}{2}$.
        \begin{enumerate}
            \item Because $\cos(x)$ is concave on $\left(0, \frac{\pi}{2}\right)$ (or $0^\circ$ to $90^\circ$), by Jensen's Inequality:
            \item $\frac{\cos(A) + \cos(B) + \cos(C)}{3} \le \cos\left(\frac{A + B + C}{3}\right) = \cos\left(\frac{\pi}{3}\right) = \frac{1}{2}$ if and only if $A = B = C = \frac{\pi}{3}$ ($60^\circ$).
            \begin{enumerate}
                \item This indicates that equality occurs only when all $3$ angles equal $60^\circ$ or $\frac{\pi}{3}$ radians (meaning $a = b = c = \sqrt{3}$).
            \end{enumerate}
        \end{enumerate}
    \end{enumerate}

    \item Compare Jensen's Inequality to the original problem:
    \begin{enumerate}
        \item Multiplying the $3$ back across the inequality sign yields $\cos(A) + \cos(B) + \cos(C) \le \frac{3}{2}$, which substitutes back into the original inequality:
        \[
            \frac{1}{\sqrt{1+a^2}} + \frac{1}{\sqrt{1+b^2}} + \frac{1}{\sqrt{1+c^2}} \le \frac{3}{2}
        \]
    \end{enumerate}
\end{enumerate}

\end{document}\documentclass[12pt]{article}
\usepackage{amsmath, amssymb, amsthm}
\usepackage{enumitem}
\usepackage[margin=1in]{geometry}

\begin{document}

\section*{Selected IMO Problem}

\begin{enumerate}
    \setcounter{enumi}{3}
    \item For positive real numbers $a, b, c$ with $a + b + c = abc$, show that
    \[
        \frac{1}{\sqrt{1+a^2}} + \frac{1}{\sqrt{1+b^2}} + \frac{1}{\sqrt{1+c^2}} \le \frac{3}{2},
    \]
    and determine when equality occurs.
\end{enumerate}

\section*{English-Language Solution}
\textbf{Proof by:} Trigonometric Substitution \& Jensen's Inequality

\begin{enumerate}
    \item Test $a, b, c$ values when they all equal one another: $(a = b = c \implies 3a = a^3) \implies a = \sqrt{3}$.
    \begin{enumerate}
        \item Plug in $a = b = c = \sqrt{3}$ to get $3\left(\frac{1}{\sqrt{4}}\right) = \frac{3}{2}$.
        \begin{enumerate}
            \item This shows that $\frac{3}{2}$ is possible when all the variables equal $\sqrt{3}$.
        \end{enumerate}
    \end{enumerate}

    \item Using the tangent identity for a triangle (conditional trigonometric identity):
    \begin{enumerate}
        \item $\tan(A) + \tan(B) + \tan(C) = \tan(A)\tan(B)\tan(C)$
        \begin{enumerate}
            \item Here, the angles must add up to $180^\circ$ ($\pi$ radians).
            \item Tangent Sum Identity: $\tan(A+B) = \frac{\tan(A) + \tan(B)}{1 - \tan(A)\tan(B)} = -\tan(C)$
            \item \textbf{Formal Proof Segment:}
            \begin{itemize}
                \item \textbf{Apply Odd Symmetry:} Since the tangent function is an odd function (meaning $-\tan(x) = \tan(-x)$ for all $x$ in its domain), we can rewrite the right-hand side of $\tan(A + B) = -\tan(C)$ as:
                \[
                    \tan(A + B) = \tan(-C)
                \]
                \item \textbf{Apply Periodicity:} The tangent function is strictly periodic with fundamental period $\pi$ (or $180^\circ$). Therefore, $\tan(\alpha) = \tan(\beta)$ if and only if $\alpha - \beta = k\pi$ for some integer $k \in \mathbb{Z}$. Applying this to $\alpha = A + B$ and $\beta = -C$ yields:
                \[
                    (A + B) - (-C) = k\pi \implies A + B + C = k\pi, \quad \text{for some } k \in \mathbb{Z}
                \]
                \item \textbf{Bound the Sum to Isolate $k$:} Since $A, B, C \in \left(0, \frac{\pi}{2}\right)$, summing these strict inequalities gives:
                \[
                    0 < A + B + C < \frac{3\pi}{2}
                \]
                Substituting $A + B + C = k\pi$ into the inequality:
                \[
                    0 < k\pi < \frac{3\pi}{2}
                \]
                Dividing through by $\pi$ gives:
                \[
                    0 < k < 1.5
                \]
                \item \textbf{Conclude:} Because $k$ must be an integer ($k \in \mathbb{Z}$), the unique solution is $k = 1$. Thus:
                \[
                    A + B + C = \pi \quad (\text{or } 180^\circ)
                \]
            \end{itemize}
        \end{enumerate}
    \end{enumerate}

    \item Set $a = \tan(A)$, $b = \tan(B)$, $c = \tan(C)$.
    \begin{enumerate}
        \item Simplify using a right triangle to get $a = \text{opp.}$ and $1 = \text{adj.}$
        \begin{enumerate}
            \item You can rewrite this to show $\cos(A) = \frac{\text{adj.}}{\text{hyp.}} = \frac{1}{\sqrt{1 + a^2}}$.
        \end{enumerate}
        \item Apply the same substitution to the rest of the expression, rewriting the goal as $\cos(A) + \cos(B) + \cos(C) \le abc $.
        \begin{enumerate}
            \item Because $\cos(x)$ is concave on $\left(0, \frac{\pi}{2}\right)$ (or $0^\circ$ to $90^\circ$), by Jensen's Inequality:
            \item $\frac{\cos(A) + \cos(B) + \cos(C)}{3} \le \cos\left(\frac{A + B + C}{3}\right) = \cos\left(\frac{\pi}{3}\right) = \frac{1}{2}$ if and only if $A = B = C = \frac{\pi}{3}$ ($60^\circ$).
            \item The first item represents the average cos(x) values of A, B, and C, which has to be less than the cos(x) of the average of the x values. This can be rewritten as cos(pi/3) since the angles must add to pi radians. 
            \item This indicates that equality occurs only when all $3$ angles equal $60^\circ$ or $\frac{\pi}{3}$ radians (meaning $a = b = c = \sqrt{3}$).
            \end{enumerate}
        \end{enumerate}
    \end{enumerate}

    \item Compare Jensen's Inequality to the original problem:
    \begin{enumerate}
        \item Multiplying the $3$ back across the inequality sign yields $\cos(A) + \cos(B) + \cos(C) \le \frac{3}{2}$, which substitutes back into the original inequality:
        \[
            \frac{1}{\sqrt{1+a^2}} + \frac{1}{\sqrt{1+b^2}} + \frac{1}{\sqrt{1+c^2}} \le \frac{3}{2}
        \]
    \end{enumerate}
\end{enumerate}

\end{document}